% Generated draft: requires mathematical and source review.
\documentclass[11pt]{article}
\usepackage[margin=25mm]{geometry}
\usepackage{fontspec}
\setmainfont{Latin Modern Roman}
\usepackage{amsmath,amssymb,mathtools,unicode-math}
\setmathfont{Latin Modern Math}
\usepackage{enumitem,needspace,xurl,hyperref}
\setlength{\parindent}{0pt}
\setlength{\parskip}{5pt}
\setlength{\emergencystretch}{3em}
\newcommand{\N}{\mathbb{N}}
\newcommand{\Z}{\mathbb{Z}}
\newcommand{\Q}{\mathbb{Q}}
\newcommand{\R}{\mathbb{R}}
\newcommand{\C}{\mathbb{C}}
\newcommand{\F}{\mathbb{F}}
\newcommand{\A}{\mathbb{A}}
\newcommand{\PP}{\mathbb{P}}
\newcommand{\E}{\mathbb{E}}
\newcommand{\eps}{\varepsilon}
\newcommand{\dd}{\,\mathrm{d}}
\begin{document}
\section*{UnsolvedMath ID: 20000224}
\textbf{Source problem code:} catalogue entry
\subsection{Definitions and mathematical statement}
The original asks whether $h(C)$ is uniformly bounded below in characteristic zero, but its extracted item omits the preceding definition of $h$. Under the proper-singularity interpretation highlighted by the report, for a reduced plane curve $C$ of degree $d$ with nonempty finite proper singular set of size $s$, the conventional quantity is
\[
 N(C)=d^2-\sum_{P\in\operatorname{Sing}C}m_P(C)^2,
 \qquad h(C)=N(C)/s.
\]
This convention excludes infinitely near points and differs from evaluating at a prescribed finite point set. The original curve class, field scope, and zero-singularity case require confirmation. In this interpretation the question is whether one constant $B\in\R$ satisfies $h(C)\geq B$ for every admissible characteristic-zero curve.

The source reports that an affirmative answer would imply Bounded Negativity for rational surfaces in characteristic zero, conventionally a lower self-intersection bound depending on the fixed surface for all reduced irreducible curves. It says $h$ is unbounded below in positive characteristic, asks what a useful replacement would be there, and reports $h(C)>-4$ for line unions with a source-era best example approximately $-3.36$. These are supplied assessments, not a current audit.

For an arrangement of $d$ distinct projective lines evaluated on its changing full proper singular locus, the research report gives an addition rule. If a new line passes through $q$ old singular points of multiplicities $m_j$ and produces $u=d-\sum_jm_j$ distinct new nodes, it claims
\[
 h'=\frac{s h+1-q-2u}{s+u}.
\]
Its arbitrary deletion formula is reported as
\[
 h'=\frac{s h+R+v-q}{s-\delta},
\]
with corrections for destroyed singularities; $R,v,q,\delta$ for that operation are not defined in the supplied summary, and no meanings are guessed.

The candidate invariant is $\Phi_t=N_t-h_0s_t$. The report claims that from any seed $h_0\leq-2$, iterative additions each meeting at most one current singular point never yield $h<h_0$. If current $h\leq-3$, any improving addition must satisfy $q\geq u+2$; for the Wiman seed it reports the exact improvement test $67q-91u>67$. All formulas require their singular-count denominators to be nonzero. This operation calculus is claimed characteristic-independent, while the main uniform-bound problem is reported open in characteristic zero.
\subsection{Short English statement}
Is the proper-singularity negativity ratio uniformly bounded below in characteristic zero? The report supplies exact line-addition/deletion rules and restrictions on certain improvement strategies, while retaining the main problem as open. The omitted definition of $h$, its singular-locus convention, and deletion corrections need source review.
\subsection{Sources}
\begin{itemize}
\item Local UnsolvedMath registry, permanent ID $20000224$, source code \texttt{catalogue entry}; created and updated 2026-08-14. Status \texttt{partially\_solved}; classification \texttt{PARTIAL-PROGRESS}; difficulty L3; no proposer or year supplied.
\item AIM Problem Lists, Algebraic geometry, workshop \emph{Degenerations in algebraic geometry}, Problem 5.1, source type PDF, \url{https://aimath.org/pastworkshops/degenalggeomproblems.pdf}. Canonical location \texttt{aim-algebraic-geometry-notes.json}, \texttt{notes[223]}, tag \texttt{problem}; statement status \texttt{exact}; original, clean, and displayed text agree. Original literature field empty.
\item Original Remark 5.2 says the question is interesting only in characteristic zero because taking all rational lines over finite fields yields unbounded negative values. A limiting construction requires varying finite fields or finite extensions, rather than infinitely many distinct all-line arrangements over a single fixed finite field. The source also motivates high-multiplicity line intersections and avoidance of simple crossings, and reports the $-4$ bound and approximate $-3.36$ example without exact citations.
\item Supplied research attempt 1: result \texttt{partial\_result}, run status \texttt{open}, rigor label \texttt{proved}, low novelty confidence. Candidate contribution: transformation formulas and an obstruction theorem for iterative additions, including the Wiman-seed criterion. No exact references, deletion-variable definitions, or proof artifacts are supplied. Explicit caveat: machine-generated research aid, not peer reviewed; independently verify before citation.
\item Dataset attribution: Ulam AI / UnsolvedMath Contributors, CC BY 4.0; original problem attribution remains to the supplied AIM workshop collection.
\end{itemize}
\end{document}
